% ============================================================
%  APPENDIX TO CHAPTER 2.  Material NOT in Rudin, imported from
%  Tao I/II, ECNU, USTC, Yu Pin.  Numbered S2.n throughout.
% ============================================================
\rappendix{2}

\begin{roadmap}[How to use this appendix]
Nothing in this appendix is Rudin's text. It is imported from the four
comparison texts and set apart at every level: the running head names the
appendix, section heads carry a reverse-video tag, item numbers are prefixed
\textbf{S}, and every statement names its source.

\medskip
\raggedright
\textbf{Why import anything at all.} Rudin's Chapter 2 is thirty pages of
topology with no payoff: every application waits for Chapter 4. He also makes
two economies that cost the reader later --- he defines \emph{limit point}
once and never names the two other formulations he goes on to use, and he
proves compactness of a cell by an inlined bisection whose key step is never
isolated as a lemma. The Chinese texts organise the same material around
\emph{what each theorem buys}, and Tao and Yu Pin carry the metric-space
theory past the point where Rudin stops.

\medskip
\textbf{What is here.}

\smallskip
\textbf{A. Cluster points.} Three inequivalent-looking definitions of ``limit
point,'' proved equivalent, and the collision between the set version and the
sequence version that Rudin's notation hides.

\medskip
\textbf{B. The completeness theorems as tools.} Nested intervals with the
corollary Rudin omits; Heine--Borel proved through it instead of by raw
bisection; Weierstrass twice; and a construction of the least upper bound
\emph{from} Cauchy completeness --- the one direction Rudin never runs,
since he takes \lub{} as an axiom.

\medskip
\textbf{C. The Heine--Borel Rudin declines.} At 2.41 he proves the
equivalences \emph{in $\R^k$ only}. The general metric-space theorem exists
--- compact $=$ complete $+$ totally bounded --- and it names exactly which
property $\R^k$ was supplying.

\medskip
\textbf{D. Two notions of smallness.} Cardinality is not the only way a set
can be small. Baire's theorem supplies the topological notion, and the two
disagree in both directions --- with Rudin's own Cantor set as one witness.

\medskip
\textbf{E. What survives without a metric.} Which of Chapter 2 is really
about distance, and which is about open sets alone.

\medskip
\textbf{How to read it.} Section A is worth reading before Chapter 3, because
the confusion it settles is the one that bites at 3.5--3.6. Section B is best
read after 2.41, as a second pass over the same theorems by different routes;
C and D after the whole chapter. E is optional and answers the question
``why a metric at all?'' The exercises are the last section.
\end{roadmap}

% ============================================================
\ssection{Supplement A --- Limit points, cluster points, and a collision}

\begin{roadmap}[Why this supplement belongs to Chapter 2]
\raggedright
\textbf{What Rudin does.} Definition 2.18(b) fixes one meaning of
\emph{limit point}: every neighbourhood of $p$ meets $E$ in a point other
than $p$. Theorem 2.20 then proves a second --- every neighbourhood meets
$E$ in \emph{infinitely many} points --- and uses it to conclude that a
finite set has no limit points.

\smallskip
\textbf{The gap.} There is a third formulation, the one that says $p$ is the
limit of a sequence of \emph{distinct} points of $E$, and Rudin never states
it. He nevertheless relies on it constantly from Chapter 3 onward: every
argument that begins ``choose a sequence in $E$ converging to $p$'' is using
it, and every such argument silently needs the chosen points to be distinct.

\smallskip
\textbf{Where the gap bites.} Twice, and both times painfully. First at
3.5--3.6, where Rudin passes between statements about \emph{sets} (2.42:
an infinite subset of a compact set has a limit point) and statements about
\emph{sequences} (3.6(b): a bounded sequence in $\R^k$ has a convergent
subsequence) as though they were the same theorem. They are not --- a
sequence has a range, and the range can be finite while the sequence has
several subsequential limits. Second at 4.2, where continuity is tested
against sequences and one must know which points those sequences may repeat.

\smallskip
\textbf{What this section supplies.} All three formulations, proved
equivalent, with the recursion that forces distinctness written out; and the
one-line example that separates the set notion from the sequence notion, so
that the translation between 2.42 and 3.6(b) can be made carefully instead
of tacitly.
\end{roadmap}

\sitem{S2.1}{Theorem}
\emph{(Three definitions of a cluster point.) Let $E \subset \R^1$ and
$\xi \in \R^1$. The following are equivalent.}
\begin{enumerate}[label=(\alph*), leftmargin=2.2em, itemsep=2pt, topsep=3pt]
  \item \emph{Every neighbourhood of $\xi$ contains \textbf{infinitely many}
        points of $E$.}
  \item \emph{Every neighbourhood of $\xi$ contains a point of $E$ other than
        $\xi$ --- Rudin's Definition 2.18(b).}
  \item \emph{There is a sequence $(x_n)$ in $E$ whose terms are
        \textbf{pairwise distinct} with $x_n \to \xi$.}
\end{enumerate}
\provenance{ECNU \cjk{聚点}{juedian} Def.\ 2, 2$'$, 2$''$ (bk.\ pp.\,151--152).
Rudin states (b) and proves (a) as 2.20; (c) he never states.}

\begin{proof}[The one implication with content]
That (a) $\Rightarrow$ (b) and (c) $\Rightarrow$ (a) are immediate. For
(b) $\Rightarrow$ (c), build the sequence recursively, and \emph{force}
distinctness by shrinking the radius past the previous point: having chosen
$x_{n-1}$, apply (b) with
\[ \varepsilon_n = \min\Bigl\{\tfrac1n,\ |\xi - x_{n-1}|\Bigr\} \]
to obtain $x_n \in E$ with $0 < |x_n - \xi| < \varepsilon_n$. Then
$x_n \to \xi$, and no point can repeat, since each is strictly nearer $\xi$
than its predecessor.
\end{proof}

\begin{unpack}[Why the recursion needs that particular radius]
The obvious recursion --- take $|x_n - \xi| < 1/n$ --- gives convergence but
not distinctness: nothing stops the same point being chosen twice. The extra
clause $|\xi - x_{n-1}|$ in the minimum makes the distances strictly
decreasing, and a strictly decreasing sequence of distances cannot repeat a
point.

This is exactly the step a reader of Rudin has to invent unaided. Theorem
2.20 gives him ``infinitely many points in every neighbourhood''; every later
argument that says ``choose a sequence in $E$ converging to $\xi$'' silently
needs those points to be distinct, and that is what the recursion supplies.
Rudin uses the fact at 2.20's corollary (a finite set has no limit points),
at 3.6, and throughout Chapter 4.
\end{unpack}

\sitem{S2.2}{Definition}
\emph{(Cluster point of a sequence.) $a$ is a cluster point of the sequence
$(x_n)$ if every neighbourhood of $a$ contains $x_n$ for \textbf{infinitely
many indices $n$} --- terms counted with repetition. Equivalently, $a$ is the
limit of some subsequence.}
\provenance{ECNU \S7.2 Def.\ 1 and the following \cjk{注}{note}
(bk.\ p.\,156).}

\begin{pitfall}[The set of terms is not the sequence]
The two notions disagree, and the cleanest witness is one line. Take
\[ x_n = \sin\frac{n\pi}{4} . \]
As a \emph{sequence} it has five cluster points: $-1$, $-\tfrac{\sqrt2}{2}$,
$0$, $\tfrac{\sqrt2}{2}$, $1$, each hit infinitely often. As a \emph{set},
its range is the five-element set
$\bigl\{-1,-\tfrac{\sqrt2}{2},0,\tfrac{\sqrt2}{2},1\bigr\}$, and a finite set
has \textbf{no} limit points at all (Rudin's corollary to 2.20).

So ``limit point'' in Chapter 2 and ``subsequential limit'' in 3.5--3.6 are
different notions that can disagree completely. Rudin uses both words without
ever flagging the difference, and 3.6(b) --- every bounded sequence in $\R^k$
has a convergent subsequence --- is \emph{not} literally 2.42 applied to the
range, because the range may be finite. The bridge is: a subsequential limit
of $(x_n)$ is either a limit point of the range, or a value taken infinitely
often. Rudin's proof of 3.6(b) quietly splits into exactly those two cases.
\end{pitfall}

\begin{center}
\begin{tikzpicture}[x=1mm, y=1mm]
  \draw[axis] (0,0) -- (104,0);
  \foreach \x/\lab in {8/{-1}, 30/{-\frac{\sqrt2}{2}}, 52/{0},
                       74/{\frac{\sqrt2}{2}}, 96/{1}}{
    \node[ptfill, minimum size=2.6pt] at (\x,0) {};
    \node[mlbl, anchor=north] at (\x,-2.2) {$\lab$};
    \foreach \k in {1,2,3,4}{
      \node[ptfill, minimum size=1.7pt] at (\x,{4.5+4*\k}) {};}
    \node[mlbl, anchor=south] at (\x,22) {$\vdots$};}
  \node[lbl, anchor=west] at (0,27.5)
    {\emph{as a sequence}: every value returns for infinitely many $n$};
  \node[lbl, anchor=west] at (0,-9.5)
    {\emph{as a set}: just these five points --- no limit point at all};
\end{tikzpicture}
\end{center}
\figcap{$x_n = \sin(n\pi/4)$ has five cluster points as a sequence and none
as a set. Multiplicity is exactly the information the range throws away ---
which is why Rudin's 2.42 (about sets) does not directly give 3.6(b) (about
sequences).}

% ============================================================
\ssection{Supplement B --- The completeness theorems as tools}

\begin{roadmap}[Why this supplement belongs to Chapter 2]
\raggedright
\textbf{What Rudin does.} He assumes the least-upper-bound property as an
axiom (1.19) and derives everything from it in one direction: 1.20
(Archimedes and density), 2.38 (nested compact sets), 2.40 ($k$-cells are
compact), 2.41 (Heine--Borel in $\R^k$), 3.11 (completeness).

\smallskip
\textbf{The gap.} A one-directional development hides the fact that these
statements are \emph{equivalent}. Chapter 1's comparison box lists six
formulations of completeness and notes that different books axiomatise
different ones; but a list of equivalent statements teaches nothing until
you have seen the implications run. Rudin runs them only outward from
\lub{}, so a reader has no way to discover that \lub{} can be recovered from
Cauchy completeness --- or what extra hypothesis that recovery needs.

\smallskip
\textbf{Where the gap bites.} At 2.40 the crucial step --- that a nested
family with vanishing diameters is eventually swallowed by any neighbourhood
of its common point --- is derived inline, in the middle of a
$k$-dimensional argument, and never named. It is needed again at 2.42, at
3.10, and in every later bisection argument, and each time it is rebuilt from
scratch.

\smallskip
\textbf{What this section supplies.} That step, isolated as a lemma with its
corollary (S2.3), after which Heine--Borel is three lines (S2.4) and
Weierstrass is four (S2.5); a proof of Bolzano--Weierstrass that uses no
bisection at all and secretly computes $\liminf$ (S2.6); and the missing
implication, a construction of $\sup$ from Cauchy completeness plus
Archimedes (S2.7), which shows precisely why the Archimedean hypothesis
cannot be dropped.
\end{roadmap}

\sitem{S2.3}{Theorem}
\emph{(Nested intervals, with its corollary.) Let $[a_{n+1},b_{n+1}] \subset
[a_n,b_n]$ for all $n$ and $b_n - a_n \to 0$. Then there is exactly one
$\xi$ with $\xi \in [a_n,b_n]$ for every $n$; and moreover, for every
$\varepsilon > 0$ there is an $N$ with}
\[ [a_n,b_n] \subset (\xi-\varepsilon,\ \xi+\varepsilon)
   \qquad (n > N). \]
\provenance{ECNU \cjk{区间套定理}{nested interval thm} 7.1 and its
\cjk{推论}{corollary} (bk.\ pp.\,150--151).}
\begin{proof}
The $a_n$ increase and every $b_m$ is an upper bound for them: if $m \le n$
then $a_n \le b_n \le b_m$, and if $m > n$ then $a_n \le a_m \le b_m$. So
$\{a_n\}$ is nonempty and bounded above, and by 1.19 it has a supremum
$\xi$, with $a_n \le \xi \le b_n$ for every $n$. If $\xi'$ also lay in every
$[a_n,b_n]$ then $|\xi-\xi'| \le b_n - a_n \to 0$, so $\xi' = \xi$.

For the corollary, given $\varepsilon>0$ choose $N$ with
$b_N - a_N < \varepsilon$. For $n > N$ the interval $[a_n,b_n]$ is contained
in $[a_N,b_N]$, has length less than $\varepsilon$, and contains $\xi$;
hence every one of its points lies within $\varepsilon$ of $\xi$.
\end{proof}

\noindent
Rudin's 2.38 gives nonempty intersection for nested compact sets and stops
there. The corollary --- eventually the whole interval, not merely the point,
sits inside the neighbourhood --- is what does the work in every application
below.\kn{The proof of the corollary is one line: $\xi \in [a_n,b_n]$ and
$b_n - a_n \to 0$, so both endpoints are within $\varepsilon$ of $\xi$ once
$b_n - a_n < \varepsilon$. Trivial to prove, easy to forget you need, and
Rudin never names it --- so at 2.40 he re-derives it inline.}

\begin{pitfall}[Closed is not decoration]
The intervals must be closed. The open intervals $(0,1/n)$ are nested with
lengths tending to $0$ and have \textbf{empty} intersection. Boundedness
matters too: the closed sets $[n,\infty)$ are nested with empty
intersection.

Rudin's 2.38 requires compactness, which packages both hypotheses at once;
seeing the two failures separately is what makes the packaging legible.
\end{pitfall}

\sitem{S2.4}{Theorem}
\emph{(Heine--Borel, via nested intervals.) Every open cover of $[a,b]$ has
a finite subcover.}
\provenance{ECNU \cjk{有限覆盖定理}{finite covering thm} 7.3
(bk.\ p.\,153).}

\begin{proofwalk}[How this proof differs from Rudin's 2.40]
  \item \textbf{Same start.} Suppose no finite subcover exists. Bisect
        $[a,b]$; at least one half again admits no finite subcover. Iterate,
        producing nested intervals with lengths $\to 0$.
  \item \textbf{The step that is now free.} Let $\xi$ be the common point.
        It lies in some cover element $(\alpha,\beta)$. By the
        \emph{corollary} to S2.3, $[a_n,b_n] \subset (\alpha,\beta)$ for
        large $n$ --- so that interval is covered by a \emph{single} member,
        contradicting the choice of intervals with no finite subcover.
  \item \textbf{What changed.} Rudin's 2.40 proves the same thing but derives
        step 2 inline, in the middle of a $k$-dimensional argument, using
        $\operatorname{diam} Q_n \to 0$. Naming the lemma first turns the
        proof into three lines and makes visible that bisection is used only
        to \emph{manufacture} a nested family --- any other method of
        manufacturing one would do.
  \item \textbf{Audit.} Closedness of $[a,b]$ enters once, in S2.3;
        boundedness enters once, to start the bisection. Both hypotheses are
        spent exactly where the pitfall above says they must be.
\end{proofwalk}

\begin{pitfall}[Where the theorem fails, made concrete]
On $(0,1)$ the cover
\[ \Bigl\{\ \Bigl(\tfrac{1}{n+1},\,1\Bigr) : n = 1,2,3,\dots \Bigr\} \]
works and no finite subfamily does: the members increase, so a finite
subfamily is its largest member, which omits everything below
$\tfrac1{N+1}$. Rudin gives the analogous fact at 2.32's examples; the point
worth keeping is the mechanism --- a totally ordered cover with no largest
element can only be finitely reduced if one member already suffices.
\end{pitfall}

\begin{center}
\begin{tikzpicture}[x=1mm, y=1mm]
  \draw[axis] (0,0) -- (104,0);
  \node[ptopen] at (6,0) {};
  \node[ptopen] at (98,0) {};
  \node[mlbl, anchor=north] at (6,-1.8) {$0$};
  \node[mlbl, anchor=north] at (98,-1.8) {$1$};
  \foreach \y/\x in {6/52, 11/33, 16/21.5, 21/14.4, 26/10.4}{
    \draw[seg] (\x,\y) -- (96,\y);
    \node[ptopen, minimum size=1.6pt] at (\x,\y) {};}
  \node[mlbl, anchor=west] at (54,6.2) {$(\tfrac12,1)$};
  \node[mlbl, anchor=west] at (35,11.2) {$(\tfrac13,1)$};
  \node[lbl, anchor=west] at (24,26) {$\cdots$ never reaching $0$};
  \draw[guide] (6,-3.4) -- (6,29);
\end{tikzpicture}
\end{center}
\figcap{The cover of $(0,1)$ that admits no finite subcover. The members are
nested upward, so any finite subfamily equals its largest member and misses a
whole interval at the left end. Compactness fails exactly at the point the
domain is missing.}

\sitem{S2.5}{Theorem}
\emph{(Weierstrass.) Every bounded infinite subset of $\R^1$ has a cluster
point.}
\provenance{ECNU \cjk{聚点定理}{cluster point thm} 7.2, with two proofs
(bk.\ pp.\,152--153).}
\begin{proof}
Let $E$ be infinite with $E \subset [-M,M]$. Bisect: at least one of the two
halves contains infinitely many points of $E$, since otherwise $E$ would be
finite. Call it $[a_1,b_1]$ and repeat, obtaining nested intervals
$[a_n,b_n]$, each containing infinitely many points of $E$, with
$b_n - a_n = 2M/2^{\,n} \to 0$.

By S2.3 there is a point $\xi$ common to all of them. Given $\varepsilon>0$,
the corollary to S2.3 puts $[a_n,b_n]$ inside $(\xi-\varepsilon,
\xi+\varepsilon)$ for all large $n$; and $[a_n,b_n]$ contains infinitely
many points of $E$. So every neighbourhood of $\xi$ contains infinitely many
points of $E$, which is S2.1(a).
\end{proof}

\begin{deeper}[Two proofs, landing on two different definitions]
\textbf{First proof (bisection).} $E \subset [-M,M]$; bisect, always keeping
a half containing infinitely many points of $E$. By S2.3 and its corollary,
every neighbourhood of the common point $\xi$ contains some $[a_n,b_n]$,
hence infinitely many points of $E$. This lands on definition (a) of S2.1.

\smallskip
\textbf{Second proof (via sequences).} Choose pairwise distinct
$x_n \in E$; the sequence is bounded, so some subsequence converges, say to
$x_0$; then every neighbourhood of $x_0$ contains infinitely many of the
$x_n$. This lands on definition (c).

\smallskip
The two proofs are the same argument in different costumes, and S2.1 is what
lets you move between them. ECNU then states the dictionary Rudin leaves
implicit: \emph{the sequence form of Bolzano--Weierstrass is the set form
applied to the set of terms, with ``infinitely many points'' read as
``infinitely many indices.''} Rudin's 2.42 and 3.6(b) sit twenty pages apart
with no remark connecting them --- and, as the pitfall at S2.2 shows, the
translation needs the care that remark would have supplied.
\end{deeper}

\sitem{S2.6}{Theorem}
\emph{(Bolzano--Weierstrass without bisection.) Every bounded sequence
$(a_n)$ of reals has a convergent subsequence.}
\provenance{USTC \cjk{定理}{Thm} 8 (bk.\ p.\,12).}

\begin{proof}[The monotone-tail proof]
Put $\alpha_n = \inf\{a_n, a_{n+1}, \dots\}$. The sequence $(\alpha_n)$ is
increasing and bounded above, hence converges; call the limit $\alpha$. Now
build indices: by the definition of infimum there is $n_1$ with
$\alpha_1 \le a_{n_1} < \alpha_1 + 1$; having chosen $n_k$, there is
$n_{k+1} > n_k$ with
\[ \alpha_{n_k} \le a_{n_{k+1}} < \alpha_{n_k} + \frac{1}{k+1} . \]
Since $\alpha_{n_k} \to \alpha$, the squeeze gives $a_{n_k} \to \alpha$.
\end{proof}

\begin{deeper}[What this proof reveals that Rudin's route hides]
Rudin reaches Bolzano--Weierstrass only through the cell machinery:
2.40 ($k$-cells are compact) $\to$ 2.42 (infinite subsets of compact sets
have limit points) $\to$ 3.6(b). That path makes the theorem look as though
it \emph{needs} compactness, and in $\R^k$ it is the efficient route.

On the line it needs nothing but the monotone convergence theorem. And the
auxiliary sequence is not arbitrary: $\alpha_n = \inf_{k \ge n} a_k$ is
precisely the quantity whose limit Rudin will call $\liminf a_n$ at 3.16. So
this proof silently establishes something stronger than it claims ---
\emph{the subsequential limit it produces is the lower limit} --- which is
Rudin's 3.17(a) arriving a chapter early and for free.

Keep the contrast: bisection halves the \emph{domain}, the monotone-tail
argument shaves the \emph{index set}. Bisection generalises to $\R^k$; the
tail argument does not, because it needs the order.
\end{deeper}

\sitem{S2.7}{Theorem}
\emph{(The least upper bound, constructed from Cauchy completeness.) Assume
$\R$ is Archimedean and that every Cauchy sequence of reals converges. Then
every nonempty set bounded above has a supremum.}
\provenance{ECNU \cjk{例}{Ex} 3 (bk.\ pp.\,154--155), the $6 \Rightarrow 1$
step of the completeness cycle.}

\begin{proofwalk}[The construction]
  \item \textbf{Manufacture near-misses.} Let $S \ne \varnothing$ be bounded
        above. For each $n$, the Archimedean property supplies an integer
        $k_n$ such that $\lambda_n = k_n/n$ is an upper bound of $S$ while
        $\lambda_n - 1/n$ is not. (Walk the multiples of $1/n$ upward until
        they first clear $S$.)
  \item \textbf{Show the near-misses are Cauchy.} Since $\lambda_m$ is an
        upper bound and $\lambda_n - 1/n$ is not, some $a' \in S$ has
        $\lambda_m \ge a' > \lambda_n - 1/n$, so
        $\lambda_n - \lambda_m < 1/n$. By symmetry
        $|\lambda_m - \lambda_n| < \max\{1/m,\,1/n\}$.
  \item \textbf{Cash in completeness.} So $\lambda_n \to \lambda$ for some
        real $\lambda$.
  \item \textbf{Identify the limit.} $\lambda$ is an upper bound, being a
        limit of upper bounds. And it is least: given $\delta>0$ choose $n$
        with $1/n < \delta/2$ and $\lambda_n > \lambda - \delta/2$; since
        $\lambda_n - 1/n$ is not an upper bound, some $a' \in S$ exceeds it,
        hence exceeds $\lambda - \delta$.
  \item \textbf{Audit.} The Archimedean property is spent once, in step 1;
        completeness once, in step 3. Neither is decorative --- see the
        pitfall.
\end{proofwalk}

\begin{deeper}[Why Rudin never does this, and why it is worth seeing]
Rudin takes the \lub{} property as an \emph{axiom} (1.19) and derives the
Archimedean property from it (1.20a). He therefore never shows the converse
direction, and a reader of Rudin alone could reasonably conclude that
``complete'' has only one meaning.

It has two, and they are not the same. \emph{Cauchy completeness alone does
not imply the \lub{} property}: there are Cauchy-complete ordered fields
without it --- non-Archimedean ones. That is why step 1 above needs the
Archimedean hypothesis, and it is exactly the gap Yu Pin's axiom list
(field + order + Archimedean + nested intervals) is built to expose.
Chapter 1's comparison box lists the equivalences; this is the proof of the
one link that Rudin's development makes structurally invisible.
\end{deeper}

% ============================================================
\ssection{Supplement C --- Completeness, and the Heine--Borel Rudin declines}

\begin{roadmap}[Why this supplement belongs to Chapter 2]
\raggedright
\textbf{What Rudin does.} Theorem 2.41 proves that for a set in $\R^k$,
three conditions coincide: compact, closed and bounded, and every infinite
subset has a limit point. He states the restriction to $\R^k$ emphatically,
and Exercise 2.16 supplies a metric space where closed and bounded fails to
give compact.

\smallskip
\textbf{The gap.} The restriction is a warning, not an explanation. A reader
finishes 2.41 knowing that ``closed and bounded'' is the wrong condition in
general, but not what the right one is --- and therefore not what property
$\R^k$ was quietly contributing.

\smallskip
\textbf{Where the gap bites.} Everywhere Rudin later works in a space that is
not $\R^k$. At 7.23 (Arzel\`a--Ascoli) he must characterise the compact
subsets of a function space, and the characterisation is exactly the general
theorem below --- pointwise boundedness and equicontinuity are total
boundedness in disguise. Chapter 11's $\mathscr L^2$ is another such space.
Without the general statement, each of these looks like a new trick.

\smallskip
\textbf{What this section supplies.} The general theorem --- compact
$=$ complete $+$ totally bounded --- with proof, so that 2.41 decodes:
\emph{closed} was supplying completeness and \emph{bounded} was supplying
total boundedness, and only the second implication is special to $\R^k$.
Also the equivalence of the two compactnesses in a metric space, which Rudin
proves only inside $\R^k$.
\end{roadmap}

\sitem{S2.8}{Definition}
\emph{A metric space $X$ is \textbf{complete} if every Cauchy sequence in $X$
converges in $X$; it is \textbf{totally bounded} if for every
$\varepsilon>0$ finitely many balls of radius $\varepsilon$ cover $X$.}
\provenance{Yu Pin \cjk{定义}{Def} 35 (p.\,76); Tao II, Ex.\ 1.5.10
(bk.\ p.\,26). Rudin defines completeness at 3.11 and total boundedness
nowhere.}

\sitem{S2.9}{Theorem}
\emph{(Heine--Borel for metric spaces.) A metric space is compact if and only
if it is complete and totally bounded.}
\provenance{Tao II, Ex.\ 1.5.10(b),(c).}
\begin{proof}
Throughout, use S2.10 to read ``compact'' as ``every sequence has a
convergent subsequence.''

\emph{Compact $\Rightarrow$ totally bounded.} Given $\varepsilon>0$, the
balls $B(x,\varepsilon)$, $x \in X$, form an open cover; a finite subcover is
a finite family of $\varepsilon$-balls covering $X$.

\emph{Compact $\Rightarrow$ complete.} Let $(x_n)$ be Cauchy. Some
subsequence converges, say $x_{n_k} \to x$. Given $\varepsilon>0$ choose $N$
with $d(x_n,x_m)<\varepsilon/2$ for $n,m \ge N$, and then $k$ with
$n_k \ge N$ and $d(x_{n_k},x)<\varepsilon/2$; for $n \ge N$ the triangle
inequality gives $d(x_n,x)<\varepsilon$. So the whole sequence converges.

\emph{Complete and totally bounded $\Rightarrow$ compact.} Let $(x_n)$ be any
sequence. Cover $X$ by finitely many balls of radius $1$; one of them
contains $x_n$ for infinitely many $n$, giving a subsequence $S_1$ lying in
a set of diameter at most $2$. Cover $X$ by finitely many balls of radius
$\tfrac12$; one contains infinitely many terms of $S_1$, giving a
subsequence $S_2 \subset S_1$ of diameter at most $1$. Continuing, we obtain
$S_1 \supset S_2 \supset \cdots$, with $S_j$ of diameter at most $2^{2-j}$.

Let $y_k$ be the $k$th term of $S_k$. For $i,j \ge k$ both $y_i$ and $y_j$
belong to $S_k$, so $d(y_i,y_j) \le 2^{2-k}$: the diagonal sequence $(y_k)$
is Cauchy, and by completeness it converges. It is a subsequence of
$(x_n)$, so $X$ is sequentially compact.
\end{proof}

\begin{proofwalk}[How the proof flows]
  \item \textbf{Choose the method.} Two properties must be produced from
        one, and then one from two. The forward direction is cheap because
        compactness is a hypothesis about \emph{every} cover and about
        \emph{every} sequence, so one simply applies it to a well-chosen
        instance. The converse is where the work is: nothing is compact yet,
        so a convergent subsequence has to be manufactured by hand.
  \item \textbf{Spend the cover hypothesis (forward, part one).} Apply
        compactness to the cover by all $\varepsilon$-balls. There is
        nothing to choose --- the cover is handed to you by the definition
        of total boundedness itself.
  \item \textbf{Spend the sequence hypothesis (forward, part two).} A Cauchy
        sequence needs only \emph{one} convergent subsequence to converge
        itself; that upgrade is the same three-line estimate as in 3.11 and
        is the whole content of ``compact implies complete.''
  \item \textbf{Manufacture nested pigeonholes (converse).} Total
        boundedness gives finitely many balls at each scale, and finitely
        many boxes holding infinitely many terms means one box holds
        infinitely many. Halving the radius each time makes the boxes shrink
        geometrically. This is pattern 3 of Chapter 1's strategy index ---
        when the hypothesis is finiteness, pigeonhole.
  \item \textbf{Take the diagonal.} No single $S_j$ converges, so no one of
        them is the answer; the diagonal belongs to every $S_k$ from the
        $k$th place on and is therefore Cauchy. The diagonal is what turns a
        \emph{nested family} of subsequences into a \emph{single}
        subsequence.
  \item \textbf{Cash in completeness.} Cauchy becomes convergent, and only
        here. Rudin performs the identical manoeuvre at 7.23 to extract a
        convergent subsequence of functions.
  \item \textbf{Audit.} Total boundedness is spent once, in step 4, to get
        the boxes; completeness once, in step 6, to convert the diagonal's
        Cauchy property into a limit. Neither is decorative: the sequence
        $(1/n)$ in $(0,1]$ is totally bounded and not complete, and the unit
        vectors of the picture below are complete and not totally bounded,
        and neither space is compact.
\end{proofwalk}

\begin{deeper}[Reading Rudin's 2.41 through this]
Now Rudin's restriction to $\R^k$ decodes. ``Closed and bounded'' is doing
two separate jobs there:

\smallskip
\emph{Closed} in $\R^k$ gives \textbf{completeness} --- a closed subset of a
complete space is complete, and $\R^k$ is complete by 3.11. \emph{Bounded}
in $\R^k$ gives \textbf{total boundedness}, because a bounded set fits in a
cell and a cell is chopped into finitely many pieces of any prescribed size.

\smallskip
The second implication is the one that fails elsewhere, and it fails for a
reason with a name: in infinite dimensions, bounded does not imply totally
bounded. Rudin's Exercise 2.16 --- the set $\{p \in \Q : 2<p^2<3\}$, closed
and bounded in $\Q$ but not compact --- diagnoses the \emph{other} failure,
incompleteness. Both diseases exist, and 2.41 is the statement that $\R^k$
has neither.
\end{deeper}

\begin{center}
\begin{tikzpicture}[x=1mm, y=1mm]
  \draw[thin] (22,14) circle (21mm);
  \node[mlbl, anchor=south] at (22,36.4) {bounded};
  \foreach \a/\b in {12/8, 30/7, 20/17, 33/19, 11/21, 24/25, 15/28.5,
                     31/28, 22/6}{
    \draw[guide] (\a,\b) circle (3.1mm);}
  \node[lbl, anchor=west] at (48,20)
    {infinitely many disjoint balls of};
  \node[lbl, anchor=west] at (48,16)
    {radius $\varepsilon$ fit inside: \emph{not}};
  \node[lbl, anchor=west] at (48,12) {totally bounded};
  \node[lbl, anchor=west] at (48,4)
    {their centres form a sequence};
  \node[lbl, anchor=west] at (48,0)
    {with no Cauchy subsequence};
\end{tikzpicture}
\end{center}
\figcap{Bounded but not totally bounded --- the failure Rudin's 2.41 rules
out by working in $\R^k$. In $\ell^1$ the unit vectors $e^{(n)}$ are at
mutual distance $2$, so the balls of radius $1$ about them are disjoint; the
set is closed, bounded, and lives in a complete space, yet no subsequence
converges. Total boundedness is precisely the hypothesis that forbids this
picture.}

\sitem{S2.10}{Theorem}
\emph{(The two compactnesses agree.) For a metric space $X$: every open cover
of $X$ has a finite subcover if and only if every sequence in $X$ has a
convergent subsequence.}
\provenance{Yu Pin, p.\,114 (both directions); Tao II, Thm 1.5.8 and
Ex.\ 1.5.11 --- Tao \emph{defines} compactness sequentially, the mirror image
of Rudin.}
\begin{proof}
\emph{Covers $\Rightarrow$ sequences.} Suppose $(x_k)$ had no convergent
subsequence. Then no point $x$ is a subsequential limit, so each $x$ has a
radius $r_x>0$ for which $B(x,r_x)$ contains $x_k$ for only finitely many
indices $k$. These balls cover $X$; a finite subcover would then account for
only finitely many indices altogether, which is absurd.

\emph{Sequences $\Rightarrow$ covers.} Let $\{U_\alpha\}$ be an open cover.
First, $X$ is totally bounded: if some $\varepsilon>0$ admitted no finite
cover by $\varepsilon$-balls, one could choose $x_1$, then $x_2$ outside
$B(x_1,\varepsilon)$, then $x_3$ outside the first two balls, and so on,
producing a sequence with mutual distances at least $\varepsilon$ and hence
no convergent subsequence.

Second, the cover has a Lebesgue number (S4.19 in the Chapter 4 appendix):
if not, then for each $n$ some $B(x_n,1/n)$ lies in no $U_\alpha$; passing
to a convergent subsequence $x_{n_k} \to x$ and choosing $\beta$ and $r$ with
$B(x,r) \subset U_\beta$, we get $B(x_{n_k},1/n_k) \subset B(x,r)$ for large
$k$ --- a contradiction.

Now cover $X$ by finitely many balls of radius $\delta$, a Lebesgue number.
Each lies inside a single $U_\alpha$, and those finitely many $U_\alpha$
cover $X$.
\end{proof}

\noindent
Rudin proves one direction inside $\R^k$ (2.41, via 2.37) and never states
the general equivalence.\kn{The forward direction is slick: if $(x_k)$ had no
convergent subsequence, every tail $\{x_k : k \ge m\}$ would be closed, so
the complements $U_m = X \setminus \{x_k : k\ge m\}$ form an open cover with
no finite subcover. The reverse direction is the one that needs a uniform
scale --- the Lebesgue number lemma, stated as S4.19 in the Chapter 4
appendix, where Rudin's 4.19 needs it too.}

% ============================================================
\ssection{Supplement D --- Two notions of smallness}

\begin{roadmap}[Why this supplement belongs to Chapter 2]
\raggedright
\textbf{What Rudin does.} The chapter opens with cardinality (2.1--2.14) and
closes with the Cantor set (2.44), an uncountable set that is, in his words,
of measure zero. Size is measured by counting.

\smallskip
\textbf{The gap.} Counting is not the only measure of size, and Rudin's own
examples show it is not the useful one for topology. $\Q$ is countable ---
small by counting --- yet meets every interval. The Cantor set is
uncountable --- large by counting --- yet contains no interval at all. The
notion that separates them is nowhere in Chapter 2, though every ingredient
for it is: closure, interior, density.

\smallskip
\textbf{Where the gap bites.} At Exercise 3.22, where Rudin poses Baire's
theorem with no indication of what it is for; at 7.18, where the existence of
a nowhere-differentiable continuous function looks like a lucky construction
rather than the typical case; and in the Chapter 4 appendix, where the
question ``can a function be continuous exactly on $\Q$?'' cannot even be
posed without it.

\smallskip
\textbf{What this section supplies.} Cantor's theorem in general form and
Schr\"oder--Bernstein, which make cardinality an order relation rather than a
collection of tricks; then Baire's theorem with proof, and the table showing
that the two notions of smallness disagree in both directions.
\end{roadmap}

\sitem{S2.11}{Theorem}
\emph{(Cantor.) For every set $X$, there is no surjection $X \to 2^X$; in
particular $X$ and its power set never have the same cardinality.}
\provenance{Tao I, Thm 8.3.1 (bk.\ p.\,195). Rudin's 2.14 is the case
$X = \N$.}
\begin{proof}
Let $f : X \to 2^X$ be any map and put
\[ A = \{\,x \in X : x \notin f(x)\,\} . \]
If $A = f(x_0)$ for some $x_0$, then $x_0 \in A$ means $x_0 \notin f(x_0) =
A$, while $x_0 \notin A$ means $x_0 \in f(x_0) = A$. Both are impossible, so
$A$ is not in the range of $f$.
\end{proof}

\noindent
The proof is Rudin's 2.14 with the decimal expansions stripped away: given
$f : X \to 2^X$, the set $A = \{x : x \notin f(x)\}$ is not $f(x)$ for any
$x$.\kn{Rudin proves $2^{\N}$ is uncountable by a diagonal argument on
sequences of $0$s and $1$s, which is this proof in a costume. Stated for
general $X$ it yields an infinite ladder of cardinalities, and with it the
question Rudin never raises --- whether anything sits strictly between
$|\N|$ and $|\R|$. That is the continuum hypothesis, and Gödel and Cohen
showed it is undecidable from the usual axioms.}

\sitem{S2.12}{Theorem}
\emph{(Schr\"oder--Bernstein.) If there are injections $A \to B$ and
$B \to A$, then there is a bijection $A \to B$.}
\provenance{Tao I, Ex.\ 8.3.3; Yu Pin p.\,38 gives a different proof by
least fixed point.}
\begin{proof}
Let $f : A \to B$ and $g : B \to A$ be injections. Put
\[ C_0 = A \setminus g(B), \qquad C_{n+1} = g\bigl(f(C_n)\bigr),
   \qquad C = \bigcup_{n\ge0} C_n , \]
and define
\[ h(a) = \begin{cases}
    f(a) & (a \in C),\\
    g^{-1}(a) & (a \notin C).
  \end{cases} \]
The second line makes sense: $a \notin C$ gives $a \notin C_0$, so
$a \in g(B)$, and $g$ is injective, so $g^{-1}(a)$ is a single point.

$h$ is injective. Both $f$ and $g^{-1}$ are, so the only risk is
$f(a) = g^{-1}(a')$ with $a \in C$, $a' \notin C$. But then $a \in C_n$ for
some $n$ and $a' = g(f(a)) \in C_{n+1} \subset C$, a contradiction.

$h$ is onto. Let $b \in B$. If $g(b) \in C$, then $g(b) \in C_{n+1}$ for some
$n$ --- it cannot lie in $C_0$ --- so $g(b) = g(f(c))$ for some $c \in C_n$,
whence $b = f(c) = h(c)$ by injectivity of $g$. If $g(b) \notin C$, then
$h(g(b)) = g^{-1}(g(b)) = b$.
\end{proof}

\noindent
This is what makes cardinality an order relation rather than a collection of
tricks: it is antisymmetry.\kn{Rudin compares cardinalities throughout
Chapter 2 --- ``at most countable,'' ``$E$ is equivalent to $\N$'' --- and
never needs to state that $|A| \le |B|$ and $|B| \le |A|$ force $|A|=|B|$,
because his cases are concrete enough to exhibit bijections directly. The
general fact is not obvious and the proof is genuinely clever.}

\sitem{S2.13}{Theorem}
\emph{(Baire.) In a complete metric space, the intersection of countably many
dense open sets is dense. Equivalently, a countable union of closed sets with
empty interior has empty interior.}
\provenance{Yu Pin \cjk{定理}{Thm} 160--161 (pp.\,296--297). Rudin gives it
as Exercise 3.22.}
\begin{proof}
Let $U_1, U_2, \dots$ be dense open sets in the complete space $X$, put
$V = \bigcap_n U_n$, and let $W$ be any nonempty open set; we must produce a
point of $V \cap W$.

Since $U_1$ is dense and open, $U_1 \cap W$ is nonempty and open, so it
contains a closed ball $\overline B(x_1,r_1)$ with $0 < r_1 < 1$. Having
chosen $\overline B(x_n,r_n)$, the set $U_{n+1} \cap B(x_n,r_n)$ is nonempty
and open, so it contains a closed ball $\overline B(x_{n+1},r_{n+1})$ with
$0 < r_{n+1} < r_n/2$.

The balls are nested, so for $m > n$ we have $x_m \in \overline B(x_n,r_n)$
and hence $d(x_m,x_n) \le r_n < 2^{-n+1}$: the centres form a Cauchy
sequence. By completeness $x_n \to x$. For each $n$, all $x_m$ with $m \ge n$
lie in the closed set $\overline B(x_n,r_n)$, so $x$ does too; therefore
$x \in U_n$ for every $n$, and $x \in \overline B(x_1,r_1) \subset W$.
\end{proof}

\begin{proofwalk}[How the proof flows]
  \item \textbf{Choose the method.} ``$V$ is dense'' means $V$ meets every
        nonempty open $W$, so the goal is to \emph{produce a point} lying in
        infinitely many sets at once. The only machinery that produces a
        point from infinitely many conditions is a nested construction
        followed by completeness --- the same shape as S2.3 and 2.38.
  \item \textbf{Use density to move, openness to get room.} At each stage
        density guarantees $U_{n+1}$ reaches into the ball we are standing
        in; openness guarantees the intersection has room for a whole closed
        ball. Both hypotheses are used at every step, and neither can be
        dropped: a dense set that is not open, like $\Q$, leaves no room.
  \item \textbf{Halve the radii.} Forcing $r_{n+1} < r_n/2$ makes the centres
        Cauchy. Any summable rate would do; halving is the cheapest.
  \item \textbf{Take \emph{closed} balls.} This is the one design choice with
        content. The limit $x$ need not lie in any open ball $B(x_n,r_n)$,
        but it must lie in the closed one, because closed sets contain their
        limits (2.27). Closing the balls is what makes the intersection
        catch the limit point.
  \item \textbf{Cash in completeness.} The centres converge, and the limit
        lies in every $\overline B(x_n,r_n) \subset U_n$, hence in $V$; and
        in $\overline B(x_1,r_1) \subset W$.
  \item \textbf{Audit.} Completeness is spent once, at step 5; density and
        openness once per stage, at step 2; closedness of the balls once, at
        step 4. Drop completeness and the theorem is false --- in $\Q$, the
        sets $\Q \setminus \{q_n\}$ are dense and open with empty
        intersection.
\end{proofwalk}

\begin{deeper}[The second notion of smallness, and how it cuts across the first]
Call a set \emph{nowhere dense} if its closure has empty interior, and
\emph{meagre} if it is a countable union of nowhere dense sets. Baire's
theorem says a complete metric space is not meagre in itself --- meagre sets
are topologically negligible.

The two notions of smallness are independent, and Chapter 2 already contains
both witnesses:

\smallskip
\begin{tabular}{@{}p{0.32\linewidth}p{0.28\linewidth}p{0.30\linewidth}@{}}
\toprule
\textbf{set} & \textbf{cardinality} & \textbf{topologically} \\
\midrule
$\Q \subset \R$        & countable --- small & dense, not meagre in itself \\
the Cantor set (2.44)  & uncountable --- large & nowhere dense --- small \\
\bottomrule
\end{tabular}

\smallskip
So the Cantor set is large by counting and negligible by topology, while
$\Q$ is the reverse. Rudin builds the Cantor set precisely to exhibit an
uncountable set of ``measure zero'' (his phrase at 2.44), which is a third
notion again --- the one Chapter 11 will develop.

Yu Pin's applications show what the theorem is for: no continuous function on
$[a,b]$ can be differentiable with unbounded derivative on every
subinterval; the indicator of $\Q$ is not a pointwise limit of continuous
functions; and ``most'' continuous functions are nowhere differentiable, the
qualitative form of Rudin's 7.18. The Chapter 4 appendix uses Baire again, at
S4.5, to show no function is continuous exactly on the rationals.
\end{deeper}

% ============================================================
\ssection{Supplement E --- What survives without a metric}

\begin{roadmap}[Why this supplement belongs to Chapter 2]
\raggedright
\textbf{What Rudin does.} He defines everything in terms of the distance
$d$: balls, limit points, open sets, compactness, connectedness. Theorem
2.24 then proves that arbitrary unions and finite intersections of open sets
are open.

\smallskip
\textbf{The gap.} 2.24 is the moment the metric stops mattering, and Rudin
does not say so. From there to the end of the chapter --- and through 4.8,
4.14 and 4.22 in the next --- the proofs use only the conclusions of 2.24,
never the distance that produced them. A reader has no way to tell which
results are really about distance and which are about open sets alone.

\smallskip
\textbf{Where the gap bites.} At 2.34 (compact implies closed), which needs a
separation hypothesis that $\R^k$ supplies for free; at the equivalence of
the two compactnesses (S2.10), which is genuinely metric and fails without
it; and at every point where a reader wonders whether uniform continuity
could be defined the same way --- it cannot, and knowing why saves a false
generalisation.

\smallskip
\textbf{What this section supplies.} The definition that takes 2.24 as an
axiom, and a table sorting Chapter 2 into what survives and what dies, with
the two counterexamples that mark the boundary.
\end{roadmap}

\sitem{S2.14}{Definition}
\emph{A \textbf{topology} on $X$ is a family $\mathcal F$ of subsets
containing $\varnothing$ and $X$ and closed under arbitrary unions and finite
intersections; its members are the open sets. $X$ is \textbf{Hausdorff} if
distinct points have disjoint neighbourhoods.}
\provenance{Tao II, \S2.5 (bk.\ pp.\,39--43); Yu Pin \cjk{定义}{Def} 78
(p.\,130).}

\begin{deeper}[Which of Rudin's Chapter 2 survives, and which dies]
Rudin's 2.24 --- arbitrary unions and finite intersections of open sets are
open --- is not a theorem in this setting but the \emph{definition}. Taking
it as the definition, one can ask which of the rest survives.

\smallskip
\begin{tabular}{@{}p{0.46\linewidth}p{0.46\linewidth}@{}}
\toprule
\textbf{Survives with no metric} & \textbf{Needs the metric} \\
\midrule
closure, interior, boundary; the relative topology (2.30); continuity as
preimages of open sets (4.8); compactness by finite subcovers; connectedness
& Cauchy sequences and completeness; boundedness; uniform continuity;
``compact $\Rightarrow$ sequentially compact'' \\
\bottomrule
\end{tabular}

\smallskip
Two of Rudin's results fail outright. \emph{Compact implies closed} (2.34)
needs Hausdorff: in the trivial topology $\{\varnothing, X\}$ every subset is
compact and only two are closed. And \emph{compact implies sequentially
compact} fails even in a Hausdorff space --- Tao's Exercise 2.5.8 constructs
a compact ordered space containing a sequence with no convergent
subsequence, which is exactly why S2.10 above had to say ``metric space.''

Uniqueness of limits is the cleanest illustration of what Hausdorff buys:
without it a sequence can converge to two different points at once. Rudin
gets uniqueness free at 3.2(b) because a metric space is automatically
Hausdorff --- take balls of radius $d(p,q)/2$.
\end{deeper}
